\scanpage{192}
ALGÈBRES de LIE SEMI-SIMPLES (9)\hfill Tribu n$^\circ$ 59

Pas de noms propres à la p. 55 ; Chevalley et Harish on parlera longuement
de dans la note historique. Dans la prop. 7 appeler les gens donnés $X'$ et $Y'$.
Au lieu du ``tout de suite'', (l.2, p. 56), dire pourquoi, (ça se propage à partir
des racines simples).

Utiliser le lemme sur les applications $R\longrightarrow$ groupe au milieu de la p. 56.
Ensuite poser $X_\alpha=u_\alpha^{-1}X'_\alpha$, (alors $t_\alpha=1$). Pas de nom à la p. 57.

\uline{\unsure{Exer6|Exerb}} - Dans l'exer. 3 donner le système des $h$, $x$, $y$ de Weyl, qui fournit des
formules sans dénominateurs. Donner la démonstration Van der Waerden, (cf.
Jacobson, bas p. 123 et p. 124).

\S4 - \uuline{Théorèmes d'existence}, (nouveau titre).

Dans ce \S\ on fait deux constructions distinctes, qui amènent presque aux
mêmes résultats. On en laissera tomber une, d'où disparition des th. 3 et 4.
Ceci amène au plan suivant :

1) Construction de $\underline a=\underline x+\underline h+\underline y$ ;

2) Définition des gros modules $E(\lambda)$ ; liberté de $\underline x$ et $\underline y$ ;

3) Construction de $\underline g$ comme quotient de $\underline a$ par les $\unsure{\Theta|\theta}_{ij}$ ; alors
$\underline g=\underline h+\sum_\lambda\underline g^\lambda$,
est de dim. finie, et est semi-simple ;

4) Définition des $W(\lambda)$ ; sont des $g$-modules ;

5) $W(\lambda)$ est de dim. finie si $\lambda$ est un pouah dominant ; conséquences pour $\underline g$.

Dans 3) la semi-simplicité de $\underline g$ se démontre ainsi : on a
$[\underline g^\alpha,\underline g^{-\alpha}]=\underline h^\alpha$,
(vrai pour $\alpha$ simple, et $W$ opère) ; alors tout idéal commutatif $\underline b$ est nul, (car
$\underline b\cap\underline g_\alpha=0$ pour toute sous-algèbre $\underline g_\alpha$ de dim. 3). A partir de là tous les cor.
concernant $\underline g$ se démontrent facilement.

\margin{191}
